\begin{table}[H]
\centering
\caption{Signed Dose-Response Estimates at Endline}
\label{tab:signed_dose_response_t1}
\footnotesize
\begin{threeparttable}
\begin{tabular}{llcccc}
\toprule
Outcome & Parameter & Estimate & Std. Error & p-value & N \\
\midrule
Regime support & Signed political dose & 0.027*** & 0.001 & <0.001 & 5,215 \\
Censorship support & Signed political dose & 0.006*** & 0.001 & <0.001 & 5,215 \\
Nationalism & Signed political dose & -0.004** & 0.001 & 0.001 & 5,215 \\
China-West thermometer gap & Signed political dose & 0.010 & 0.022 & 0.635 & 5,215 \\
Trust in foreign media & Signed political dose & 0.011*** & 0.001 & <0.001 & 5,215 \\
\bottomrule
\end{tabular}
\begin{tablenotes}[flushleft]
\footnotesize
\item Note: Entries report OLS dose-response estimates on endline respondents. `Signed political dose' equals the number of scheduled political China stories, coded positively for pro-China treatments and negatively for anti-China treatments. All specifications include an apolitical-China indicator, the baseline value of the dependent variable, and randomization-block fixed effects. In the endline sample, realized political openings equal scheduled political dose exactly, so an IV specification adds no identifying variation beyond this design-based dose regression. Heteroskedasticity-robust HC2 standard errors are reported in the table.
\end{tablenotes}
\end{threeparttable}
\end{table}
